% Design/research fragment. Included by the main cost report.
% No numerical cost conclusion is asserted by this fragment.
\section{Construction choice and what the estimates establish}

The design target is a read service for $N=2^{28}$ key--value pairs, each
containing a 32-byte key and a 32-byte value, at 10,000 logical requests per
second. The security target is ten parties with at most four malicious
corruptions and security with abort. Mean response time should remain below
five seconds, with lower latency preferred. The raw database is 16\,GiB;
cuckoo slack, tree buckets, shares, authentication, recursive maps, and
independent service replicas add to this substantially. Logical reads still
modify the ORAM's position maps and placement state. Application writes and
their replication protocol are outside this estimate.

The proposed starting point is two-choice cuckoo hashing over recursive
Circuit OPRAM, with both candidates probed on every lookup. Two separate
cuckoo tables permit two independent OPRAMs to run in parallel; this modestly
reduces individual pool sizes without requiring secret-dependent shard
routing. The model uses single-slot hashing with explicit slack and a scanned
overflow stash. Hash-table and ORAM stash capacities remain parameters that
need a failure analysis. If clients supply secret shares of indices computed
using public hash functions, the interface must specify which input
consistency properties are checked. A keyed hash evaluated inside MPC is a
different, additional cost; it has not been selected and is not implicitly
free.

\paragraph{Circuit OPRAM is more than a batch of Circuit ORAM paths.}
For $b$ probes into a binary tree with leaf depth $H$, let
$r=\lceil\log_2(2b)\rceil$. The parallel construction keeps disjoint subtrees
below depth $r$ and represents the upper tree and stash in a pool. Its
algorithms include private conflict resolution, pool retrieval, eviction
candidate selection, and pool cleanup. The deterministic eviction paths
occupy distinct subtrees. These are essential costs, not an optional
implementation detail.\footnote{Chan and Shi, \emph{Circuit OPRAM: Unifying
Statistically and Computationally Secure ORAMs and OPRAMs},
\url{https://eprint.iacr.org/2016/1084}.}

For accounting, the reference template uses
\[
 P=Z(2^r-1)+R_{\mathrm{stash}},\qquad h=H-r+1,
\]
where $Z$ is bucket capacity and $R_{\mathrm{stash}}$ is an exposed slack
parameter. A conservative temporary allocation allows $P+3b$ records while
fetched blocks and returned eviction candidates are processed. The pool
replaces the upper tree buckets; their storage must not be counted twice.
Recursion stops before an additional tree would have fewer than $2b$ leaves.
This capacity accounting does not certify the chosen concrete stash slack.

\paragraph{Reference compilation versus the best achievable system.}
The program constructs an explicit operation ledger for a plausible
Circuit-OPRAM-inspired compilation. It includes pool operations, metadata
conversion, record movement, duplicate handling, and eviction planning. Some
scan and planning coefficients are deliberately exposed estimates. It is
neither a complete circuit implementation nor a proven upper bound on every
implementation detail. Large reduction counts therefore demonstrate the
expense of this specified compilation; they are not a lower bound on MPC
lookup, a proof that the throughput goal is infeasible, or a claim to have
found the globally optimal construction. AWS prices translate these
conditional workloads into dollars; they do not validate the workload or
processor-throughput assumptions.

The accounting distinguishes scalar products, private inner-product outputs,
local arithmetic, random bits, openings, memory traffic, and dependent
communication phases. In an appropriate Shamir backend, an inner product
can communicate per output rather than per scalar term. Its local
computation still grows with its length. This distinction is particularly
valuable for retrieval and flat position-map bases.

\section{Optimizations with a concrete implementation path}

\paragraph{Use private shuffles to reduce repeated payload movement.}
The optimized reference sort augments each sorting key with a secret unique
ordinal, securely shuffles records, and then opens comparison outcomes.
Subsequent sorting moves are public local operations. The unique ordinal is
essential: otherwise duplicate-key comparison patterns can disclose
multiplicity. Vectorizing comparisons across quicksort recursion levels
exposes parallelism, although its height and latency remain random.
This shuffle-then-sort approach is described in the Orq implementation.
\footnote{\emph{Orq: Complex Analytics on Private Data with Strong Security
Guarantees}, SOSP 2025,
\url{https://cs-people.bu.edu/vkalavri/assets/pdf/orq-sosp25.pdf}.}

A generic realization has five fixed parties each choose a uniform private
permutation and input checked switching-network controls. Under static
corruption of at most four parties, at least one permutation remains hidden.
Uniform switch bits alone would not produce uniform permutations. Restoring
request order requires both undoing public sorting moves and applying the
secret networks in reverse; the latter is charged. This is a useful
reproducible baseline, but five sequential networks are expensive.

For a more ambitious implementation, benchmark a Shamir-specific permutation
primitive. The 2026 Shuffle2 work uses inner products and preprocessing to
obtain linear online shuffle communication in the malicious honest-majority
setting. Its published implementation measures semi-honest execution; the
paper's malicious overhead factors are estimates. Its offline cost also
remains substantial. Consequently, it motivates an optimization experiment,
not an unqualified replacement throughput assumption.\footnote{Gao et al.,
\emph{SLIDE: Shuffle Shamir Secret Shares Uniformly with Linear Online
Communication and Guaranteed Output Delivery}, especially Sections 5 and
Appendix C.4, \url{https://arxiv.org/html/2608.27129v1}.
Implementation: \url{https://github.com/GJCPP/ShamirShuffle}.}

\paragraph{Specialize pool operations rather than sorting everything alike.}
Pool cleanup is stable compaction on validity bits. Local prefix sums and
scalar products can compute a unique secret destination permutation, leaving
private routing as the expensive primitive; long-key comparisons are
unnecessary. Maintaining the pool in logical-address order also permits a
derived finite-difference gather: attach differences of adjacent payloads to
pool records, insert zero increments for requests, merge by address, and use
local prefix sums to recover candidate values. Secret equality still decides
whether an exact match exists. Candidate-selection routing must subsequently
be undone and stable compaction must preserve the sorted-pool invariant.
This optimization needs an implemented invariant check and security argument;
it is not assumed by the reference scan coefficients.

\paragraph{Choose metadata representation separately from payload encoding.}
Moving one field word and comparing a secret integer have different costs.
Repeatedly converting compact arithmetic metadata to bits is expensive;
retaining every address and leaf bit as an arithmetic share expands both
storage and routing. A practical comparison should include compact metadata,
retained bit metadata, and a mixed representation with narrow comparisons
and explicit conversion costs. Retained bits must be charged whenever their
record moves. Leaf compatibility against a public eviction path can reuse
prefix-match products; target-depth comparisons use roughly $\log_2 H$ bits,
rather than the full key width.

The published depth analysis of Circuit ORAM provides a starting point for
parallel eviction planning. Its word-depth bound is not directly an MPC-round
bound: comparison, selection, and authentication phases remain to be
compiled.\footnote{Chan, Chung, and Shi, \emph{On the Depth of Oblivious
Parallel RAM}, \url{https://eprint.iacr.org/2017/861}.}

\paragraph{Tune recursive packing and the base together.}
Packing $F$ parent labels per recursive block reduces record counts but
increases record width, private subblock selection, and routing cost.
Moderate powers of two, for example $F=8$ through $128$, are useful search
points rather than a predetermined optimum. The binary eviction tree can
remain unchanged. Several requests may update different offsets of the same
map block: group those updates privately and commit the block once.
Independent whole-block writes would overwrite each other.

The base contains one scalar label per record in the last instantiated tree.
It is not smaller by another factor of $F$ unless another recursive tree has
been paid for. With $n=2^d$ base labels, a divide-and-conquer one-hot decoder
has the derived multiplication recurrence
\[
 T(d)=T(\lfloor d/2\rfloor)+T(\lceil d/2\rceil)+2^d,
 \qquad T(0)=T(1)=0,
\]
and depth $\lceil\log_2 d\rceil$. For $n=4096$, this is 4,272 products per
request and four product layers, before validity gating and read/write
operations. After duplicate resolution, dot-product retrieval and aggregated
updates need about $b+n$ reduction outputs, while local work remains
proportional to $bn$. Compare this against a complete additional recursive
level. The reference program uses the divide-and-conquer decoder, adding one
product and a conservative extra phase for a secret representative-valid bit.

\paragraph{Optimize scheduling after retaining dependencies.}
At one level, pool fetch and subtree fetch can overlap once old labels are
available. Maintenance on an already accessed recursive tree can overlap
fetching from the next tree if remapping labels have been sampled beforehand
and state updates are coordinated. Recursion still has a base-to-data fetch
dependency, and the next batch must not observe an uncommitted state. A sum
of all phase counts is a conservative serial schedule; replacing it by the
maximum level depth would be unjustified. Centrally forming batches before
dispatching them to independent replicas avoids multiplying formation delay
by the number of replicas. At 10,000 requests/s, batch size $Q$ alone adds
approximately $(Q-1)/(2\cdot10{,}000)$ seconds of mean formation delay.

\section{Validation needed before using a deployment budget}

The next experiment should implement and measure one complete recursive
batch, including duplicates and misses, on the selected malicious backend.
Separately measure comparison/conversion, shuffle and inverse routing, pool
cleanup, subblock updates, eviction, and preprocessing replenishment. Then
measure the composed dependency schedule and memory footprint, rather than
inferring end-to-end performance from multiplication throughput alone.

Failure parameters must cover the intended deployment lifetime. One year at
10,000 requests/s contains about $2^{38.2}$ logical requests, before counting
two probes, recursive levels, replicas, and checking events. If $A$ events
each have failure probability at most $\epsilon$, a simple lifetime bound is
$A\epsilon$. For example, an annual target of $2^{-40}$ with $2^{42}$
relevant events would require per-event probability at most $2^{-82}$.
The ORAM stash, cuckoo overflow, statistical checks, and any randomized
sorting limits need their own justified bounds. A simulated absence of
overflow is useful testing evidence but does not establish such a bound.

Finally, capacity should include authentication and preprocessing, local
routing between machines of one party, batch queues, and the full memory of
each independently evolving replica. Cross-availability-zone placement
changes both communication charges and latency. The resulting estimate is a
proposal-planning scenario until an end-to-end malicious implementation has
been calibrated on the intended topology.
